\chapter{Captura y Análisis de Requisitos}
\label{sec:propuesta}


\begin{comment}
En esta sección se describe la propuesta del perfil de tesis/proyecto de grado. Dicha propuesta debe estar basada en el marco teórico, a fin de facilitar la comprensión del lector acerca de la innovación tecnológica que el candidato pretende llevar a cabo durante el desarrollo de su tesis/proyecto de grado.
La propuesta de tesis/proyecto de grado debe hacer uso de diversos recursos, tales como: ecuaciones, diseños CAD, circuitos de diferente índole, tablas, figuras, entre otras que permitan describir con efectividad los objetivos que desea alcanzar en caso de ser aprobado su perfil.
Este apartado o sección, por tanto, es de gran importancia y debe ser detallado con la objetividad y el profesionalismo que se espera de un candidato al grado de ingeniero.
\end{comment}

%---------------------------------------------------------------------------------------------
%---------------------------------------------------------------------------------------------
\section{Propuesta de Proyecto de Grado}


%---------------------------------------------------------------------------------------------
\subsection{Análisis de Requisitos}


\section{Análisis de Costos}

%---------------------------------------------------------------------------------------------
%---------------------------------------------------------------------------------------------
\section{Cronograma de actividades}
Como se puede apreciar en la Tabla \ref{tab:cronograma_proyecto}, asi como en el diagrama Gantt del apéndice \ref{sec:cronograma}, se presenta un cronograma que detalla las fases y actividades del proyecto, que se desarrolla desde el 1 de febrero del 2024 hasta el 29 de septiembre de 2024. El mismo, siguiendo la metodología PXP se divide en tres fases principales mas la adición de una cuarta a fines de evaluar el desempeño en actividad del producto. Estas etapas son: Definición de Requisitos, Planificación, Iteraciones de Desarrollo, y Cierre. Cada fase incluye entregables específicos y actividades que se ejecutarán en un periodo determinado para asegurar la implementación exitosa del sistema.

\begin{table}[!htb]
    \caption{Cronograma del proyecto: objetivos 1 a 3.}
    \label{tab:cronograma_proyecto}
    \centering
    \footnotesize
    \setstretch{1}
    \renewcommand{\arraystretch}{1.08}
    \begin{tabular}{|p{0.62\textwidth}|c|c|c|}
        \hline
        \rowcolor{lightgray}
        \textbf{Objetivo y actividad} & \textbf{Inicio} & \textbf{Fin} & \textbf{Días} \\
        \hline
        \rowcolor{lightgray!50}
        \textbf{Objetivo 1: requerimientos} & 01/07/2026 & 12/11/2026 & 135 \\
        \hline
        Revisión de literatura y definición del tema & 01/07/2026 & 01/09/2026 & 63 \\
        \hline
        Relevamiento en la Panadería San Miguel & 02/09/2026 & 05/09/2026 & 4 \\
        \hline
        Procesamiento de cantidades y tiempos & 06/09/2026 & 16/09/2026 & 11 \\
        \hline
        Redacción del capítulo 1 & 06/09/2026 & 05/10/2026 & 30 \\
        \hline
        Actualización del capítulo 1 con los ingredientes secos & 21/10/2026 & 28/10/2026 & 8 \\
        \hline
        Trazabilidad y jornada ordinaria pendiente & 04/11/2026 & 08/11/2026 & 5 \\
        \hline
        Lista de requerimientos & 07/11/2026 & 11/11/2026 & 5 \\
        \hline
        Redacción del objetivo 1 & 11/11/2026 & 12/11/2026 & 2 \\
        \hline
        \rowcolor{lightgray!50}
        \textbf{Objetivo 2: tolerancias} & 09/11/2026 & 14/11/2026 & 6 \\
        \hline
        Aplicación de OIML R 61 a las masas de la receta & 09/11/2026 & 11/11/2026 & 3 \\
        \hline
        Propuesta de tolerancias por insumo & 11/11/2026 & 13/11/2026 & 3 \\
        \hline
        Redacción del objetivo 2 & 13/11/2026 & 14/11/2026 & 2 \\
        \hline
        \rowcolor{lightgray!50}
        \textbf{Objetivo 3: diseño mecánico y de medición} & 14/10/2026 & 15/01/2027 & 94 \\
        \hline
        Redacción del marco teórico & 14/10/2026 & 14/11/2026 & 32 \\
        \hline
        Selección del mecanismo y cotización & 04/11/2026 & 13/11/2026 & 10 \\
        \hline
        Medición de propiedades de los insumos & 09/11/2026 & 12/11/2026 & 4 \\
        \hline
        Dimensionamiento y selección de componentes & 10/11/2026 & 16/11/2026 & 7 \\
        \hline
        CAD preliminar & 14/11/2026 & 18/11/2026 & 5 \\
        \hline
        Redacción de la propuesta y del análisis económico & 16/11/2026 & 19/11/2026 & 4 \\
        \hline
        Entrega del perfil & 20/11/2026 & 20/11/2026 & 1 \\
        \hline
        Preparación de la defensa y correcciones & 21/11/2026 & 14/12/2026 & 24 \\
        \hline
        Defensa del perfil & 07/12/2026 & 07/12/2026 & 1 \\
        \hline
        CAD de detalle y planos & 15/12/2026 & 15/01/2027 & 32 \\
        \hline
    \end{tabular}

    \vspace{2pt}
    {\footnotesize Fuente: Elaboración propia. Días calendario, con actividades en paralelo.}
\end{table}

\begin{table}[!htb]
    \caption{Cronograma del proyecto: objetivos 4 a 7 y cierre.}
    \label{tab:cronograma_proyecto_2}
    \centering
    \footnotesize
    \setstretch{1}
    \renewcommand{\arraystretch}{1.08}
    \begin{tabular}{|p{0.62\textwidth}|c|c|c|}
        \hline
        \rowcolor{lightgray}
        \textbf{Objetivo y actividad} & \textbf{Inicio} & \textbf{Fin} & \textbf{Días} \\
        \hline
        \rowcolor{lightgray!50}
        \textbf{Objetivo 4: control} & 15/12/2026 & 12/02/2027 & 60 \\
        \hline
        Selección del controlador & 15/12/2026 & 24/12/2026 & 10 \\
        \hline
        Modelo dinámico de la alimentación y el pesaje & 26/12/2026 & 15/01/2027 & 21 \\
        \hline
        Diseño y simulación del controlador & 16/01/2027 & 05/02/2027 & 21 \\
        \hline
        Redacción del objetivo 4 & 06/02/2027 & 12/02/2027 & 7 \\
        \hline
        \rowcolor{lightgray!50}
        \textbf{Objetivo 5: implementación} & 16/12/2026 & 14/03/2027 & 89 \\
        \hline
        Compra de componentes, con margen de importación & 16/12/2026 & 31/01/2027 & 47 \\
        \hline
        Fabricación de tolvas y estructura & 18/01/2027 & 14/02/2027 & 28 \\
        \hline
        Electrónica, firmware e interfaz & 25/01/2027 & 21/02/2027 & 28 \\
        \hline
        Montaje, integración y puesta a punto & 15/02/2027 & 07/03/2027 & 21 \\
        \hline
        Redacción del objetivo 5 & 08/03/2027 & 14/03/2027 & 7 \\
        \hline
        \rowcolor{lightgray!50}
        \textbf{Objetivo 6: evaluación} & 15/02/2027 & 14/04/2027 & 59 \\
        \hline
        Protocolo de pruebas & 15/02/2027 & 21/02/2027 & 7 \\
        \hline
        Pruebas funcionales y ajustes & 08/03/2027 & 21/03/2027 & 14 \\
        \hline
        Comparación con el procedimiento manual & 22/03/2027 & 04/04/2027 & 14 \\
        \hline
        Análisis y redacción del marco analítico & 05/04/2027 & 14/04/2027 & 10 \\
        \hline
        \rowcolor{lightgray!50}
        \textbf{Objetivo 7: costos} & 22/03/2027 & 09/04/2027 & 19 \\
        \hline
        Desglose del costo de fabricación & 22/03/2027 & 04/04/2027 & 14 \\
        \hline
        Evaluación del retorno de la inversión & 05/04/2027 & 09/04/2027 & 5 \\
        \hline
        \rowcolor{lightgray!50}
        \textbf{Documento final} & 10/04/2027 & 30/04/2027 & 21 \\
        \hline
        Conclusiones y artículo en formato IEEE & 10/04/2027 & 18/04/2027 & 9 \\
        \hline
        Revisión del tutor y correcciones & 12/04/2027 & 28/04/2027 & 17 \\
        \hline
        Entrega del documento final & 30/04/2027 & 30/04/2027 & 1 \\
        \hline
    \end{tabular}

    \vspace{2pt}
    {\footnotesize Fuente: Elaboración propia. Días calendario, con actividades en paralelo.}
\end{table}


La primera fase, Definición de Requisitos, abarca desde el 1 de febrero hasta el 10 de mayo de 2024, e incluye entregas clave como el marco referencial, avances y la defensa del perfil del proyecto. La segunda fase, Planificación, se extiende desde el 15 de marzo hasta el 30 de junio de 2024, y cubre actividades como la entrega de avances adicionales, la entrega completa del perfil, la defensa final del perfil, el análisis de procesos y el diseño de la arquitectura de alto nivel. La tercera fase, Iteraciones de Desarrollo, se lleva a cabo del 1 de julio al 1 de septiembre de 2024, con tres iteraciones de desarrollo, cada una de tres semanas. Finalmente, la fase de Cierre y Evaluación Final, del 2 de septiembre al 29 de septiembre de 2024, incluye la preparación de documentación, capacitación de usuarios, monitoreo inicial y ajustes, y pruebas con usuarios.

En conclusión, el cronograma establece un plan detallado para la implementación del sistema automatizado de gestión de PGSST, asegurando que cada fase del proyecto se ejecute de manera organizada y meticulosa. Este enfoque estructurado permite abordar de manera efectiva los objetivos del proyecto, garantizando que todas las actividades necesarias se completen dentro de los plazos establecidos, lo que facilita un desarrollo coherente y exitoso del sistema que cumplirá con los estándares y requisitos normativos en Bolivia.

%---------------------------------------------------------------------------------------------
%---------------------------------------------------------------------------------------------
\section{Indice tentativo}
En esta sección se describen de forma detallada todos los capítulos a ser desarrollados en el proyecto de grado. En la introducción, se establecen los antecedentes de la problemática. Se plantea y define el problema a investigar, se delinean los objetivos generales y específicos, y se justifica la relevancia del mismo desde perspectivas legal, tecnológica y social. Además, se delimitan los límites y alcances de la investigación, proporcionando un marco contextual sólido para el desarrollo del proyecto.

El desarrollo del trabajo se fundamenta en el marco teórico y legal, donde se exploran conceptos fundamentales de salud y seguridad ocupacional, así como la legislación aplicable a la materia. Se abordan también metodologías de desarrollo de software, análisis y evaluación de software, y se profundiza en el campo de la inteligencia artificial, proporcionando un contexto amplio y diverso que sustenta la investigación los objetivos del proyecto.

Tras el marco practico en el que se resuelve el problema. Finalmente, se presentan las conclusiones derivadas del proyecto, se ofrecen recomendaciones pertinentes basadas en los hallazgos obtenidos, y se esbozan posibles líneas de trabajo. Este índice proporciona una guía clara y estructurada del contenido que se desarrollará a lo largo del trabajo, asegurando una presentación coherente y completa del proyecto.

\begin{enumerate}[label*=\arabic*.]
    \item \textbf{Introducción}
    \begin{enumerate}[label*=\arabic*.]
        \item \textbf{Antecedentes}
        \item \textbf{Planteamiento del Problema}
        \item \textbf{Definición del Problema}
        \item \textbf{Objetivos}
        \begin{enumerate}[label*=\arabic*.]
            \item Objetivo general
            \item Objetivos específicos
        \end{enumerate}
        \item \textbf{Justificación}
        \begin{enumerate}[label*=\arabic*.]
            \item Justificación legal
            \item Justificación tecnológica
            \item Justificación social
        \end{enumerate}
        \item \textbf{Delimitación}
        \begin{enumerate}[label*=\arabic*.]
            \item Límites
            \item Alcances
        \end{enumerate}
    \end{enumerate}
    \item \textbf{Marco teórico y legal}
    \begin{enumerate}[label*=\arabic*.]
        \item \textbf{Salud y Seguridad Ocupacional}
        \begin{enumerate}[label*=\arabic*.]
            \item Conceptos Fundamentales
            \begin{enumerate}[label*=\arabic*.]
                \item Seguridad Ocupacional
                \item Salud Ocupacional
                \item Higiene Ocupacional
                \item Enfermedad Ocupacional
                \item Sistema de Gestión de Riesgos Ocupacionales
            \end{enumerate}
            \item Accidentes e Incidentes Ocupacionales
            \item Factores de accidentes
            \begin{enumerate}[label*=\arabic*.]
                \item Actos Inseguros
                \item Condición Insegura
            \end{enumerate}
            \item Ingeniería de Seguridad
            \item Práctica de Seguridad
            \item Investigación de accidentes
            \item Peligro
            \item Riesgo
            \begin{enumerate}[label*=\arabic*.]
                \item Riesgo Físico
                \item Riesgo Mecánico
                \item Riesgo Químico
                \item Riesgo Biológico
                \item Riesgo Psicosocial
                \item Riesgo Ergonómico
                \item Riesgo Aceptable
            \end{enumerate}
            \item Condiciones de Trabajo
            \begin{enumerate}[label*=\arabic*.]
                \item Iluminación
                \item Estrés Térmico
                \item Sonometría
                \item Ventilación
                \item Señalización
                \item Ergonomía
                \item Equipo de Protección Personal (EPP)
                \item Comité Mixto
                \item Inspección de Salud y Seguridad en el trabajo
            \end{enumerate}
        \end{enumerate}
        \item \textbf{Legislación Aplicable a la SySO}
        \begin{enumerate}[label*=\arabic*.]
            \item Contexto Histórico
            \item Normativa obligatoria
            \begin{enumerate}[label*=\arabic*.]
                \item Constitución Política del Estado
                \item Ley General del Trabajo
                \item Reglamento de la Ley General de Trabajo
                \item Ley General de Higiene, Seguridad Ocupacional y Bienestar
                \item Disposiciones Complementarias
            \end{enumerate}
            \item Normativa voluntaria
            \begin{enumerate}[label*=\arabic*.]
                \item NB/ISO 62005:2005
                \item NB/ISO 55001:2005
                \item NB/ISO 51002:2012
                \item NB/ISO 7243:2018
                \item NB/ISO 45001:2018
                \item NB/ISO 51001:2022
                \item NB/ISO 58005:2022
                \item NB/ISO 11226:2022
            \end{enumerate}
        \end{enumerate}
        \item \textbf{Metodologías de Desarrollo de Software}
        \begin{enumerate}[label*=\arabic*.]
            \item Programación Extrema
            \item Programación Extrema Personal
            \begin{enumerate}[label*=\arabic*.]
                \item Técnica MoSCoW
            \end{enumerate}
        \end{enumerate}
        \item \textbf{Análisis y Evaluación de Software}
        \begin{enumerate}[label*=\arabic*.]
            \item Estándar ISO 25000:2014
            \item Método Objetivo-Pregunta-Métrica
        \end{enumerate}
        \item \textbf{Inteligencia Artificial}
        \begin{enumerate}[label*=\arabic*.]
            \item Aprendizaje
            \item Redes Neuronales
            \item Visión Computacional
            \item Procesamiento de Lenguaje Natural (PLN)
            \item Modelos de Lenguaje Multimodales
            \begin{enumerate}[label*=\arabic*.]
                \item Transformadores
            \end{enumerate}
            \item Evaluación de Modelos
            \item Generación de Lenguaje Natural
            \item Generación de Descripciones de Imágenes
        \end{enumerate}
    \end{enumerate}
    \item \textbf{Marco Practico}
        \begin{enumerate}[label*=\arabic*.]
            \item Revisión de Requisitos
            \item Casos de Uso
            \item Selección de Tecnologías
            \item Arquitectura del Sistema  
            \item Implementación del sistema
            \item Evaluación de la calidad del software
        \end{enumerate}
    \item \textbf{Marco Analítico}
        \begin{enumerate}[label*=\arabic*.]
            \item Resultados y Discusión
            \item Análisis de Costos
                \begin{enumerate}[label*=\arabic*.]
                    \item Costos de Desarrollo
                    \item Costos de Infraestructura
                    \begin{enumerate}[label*=\arabic*.]
                        \item Servidores y Almacenamiento
                        \item Red y Seguridad
                        \item Costos Totales de Infraestructura
                    \end{enumerate}
                    \item Gastos Legales
                    \item Inversión Total Inicial
                \end{enumerate}
            \item Evaluación de Proyecto de Inversión
            \begin{enumerate}[label*=\arabic*.]
                \item Comparación Costo Beneficio
                \item Valor Actual Neto
                \item Tasa Interna de Retorno
                \item Conclusión
            \end{enumerate}
            \item Análisis de Competidores
            \item Conclusión del capítulo
        \end{enumerate}
    \item \textbf{Marco Conclusivo}
        \begin{enumerate}[label*=\arabic*.]
            \item Conclusiones
            \item Recomendaciones
            \item Trabajos Futuros
        \end{enumerate}
    \item \textbf{Bibliografía, Apéndices y Anexos}
\end{enumerate}